%=======================================================================
\chapter{Linear Elasticity}
%=======================================================================

Linear elasticity is the theory that carries almost all of engineering
practice, and it is obtained from Chapter 11 by one approximation:
the displacement gradient is small. The chapter works out what that
approximation does to each object built in the preceding eleven
chapters, and it takes some care over one point that most treatments
pass over, the shifter.

For simplicity the theory is developed in its purely mechanical form,
the free energy being independent of temperature and the temperature
itself uniform and constant throughout the body. By \eqr{11.51} and
Problem 14 of Chapter 11, the thermal energy balance is then satisfied
identically provided the internal heating supply vanishes: with
$\grad\theta=\bze$ the flux vanishes by \eqr{11.55}, and with
$\dot\theta=0$ and $W$ independent of $\theta$ the coupling term
$\theta(W_{\theta})^{\textstyle\cdot}$ vanishes too.

%=======================================================================
\section{Kinematics}\label{sec:ch12kin}
%=======================================================================

The premise of the theory is that the displacement gradient $\bH$ of
\eqr{7.22} is small everywhere in the body. Taking $\bH$ continuous in
$\bX$, this is
\begin{equation}
   \varepsilon=\max_{\bX\in\kappa}|\bH(\bX,t)|\ll1 ,
   \tag{12.1}\label{eq:12.1}
\end{equation}
where $|\bH|^{2}=\tr(\bH\bH^{t})=\tr(\bH^{t}\bH)$, so that
$|\bH|=\sqrt{H_{iA}H_{iA}}$ in Cartesian components. Since a sum of
squares dominates each of its terms, $|H_{iA}|\ll1$ for every pair
$iA$.

\begin{remark}[the letter $\varepsilon$ in this chapter]%
\label{rem:ch12eps}
Three near-identical symbols are about to appear, and the book uses all
three. The scalar $\varepsilon$ of \eqr{12.1} is the small parameter.
The referential infinitesimal strain of \eqr{12.8} is
$\vek{\epsilon}$, the closed epsilon. The spatial infinitesimal strain
of \eqr{12.17} is $\vek{\varepsilon}$, the open one. Both strains carry
the under-tilde that marks a tensor, so the scalar is distinguishable
at sight; the two tensors are told apart by their shape, and the
shifter is what converts one into the other. The internal energy, which
also wore the letter $\varepsilon$ in Chapter 6, does not appear here at
all, the theory being purely mechanical.
\end{remark}

\subsection*{Strain}

By \eqr{7.22}, $\bF=\shift+\bH$, so the right Cauchy--Green tensor is
\begin{equation}
   \bC=\bF^{t}\bF=\hI+\shift^{t}\bH+\bH^{t}\shift+\bH^{t}\bH ,
   \tag{12.2}\label{eq:12.2}
\end{equation}
the first term because $\shift^{t}\shift=\hI$, the shifter being an
isometry. The three remaining terms are all referential, but the middle
two are built from a two-point tensor, and it is worth introducing the
object that removes that awkwardness. Set
\begin{equation}
   \bC=\hI+\bG+\bG^{t}+\bG^{t}\bG,\qquad\text{where}\quad
   \bG=\shift^{t}\bH
   \tag{12.3}\label{eq:12.3}
\end{equation}
is a referential tensor with components $G_{AB}=\dd_{iA}H_{iB}$. The
identification of the last term uses
$\bG^{t}\bG=\bH^{t}\shift\shift^{t}\bH=\bH^{t}\I\bH=\bH^{t}\bH$, and of
the second $\bG^{t}=(\shift^{t}\bH)^{t}=\bH^{t}\shift$.

What $\bG$ is becomes clear on applying it to a material line element:
\begin{equation}
   \bG\,\dl\bX=\shift^{t}\bH\,\dl\bX=\shift^{t}\,\dl\bu=\dl\bw,
   \qquad\text{where}\quad \bw=\shift^{t}\bu
   \tag{12.4}\label{eq:12.4}
\end{equation}
is the \emph{referential displacement}, $\bu$ being the spatial
displacement of \eqr{2.77}, and where the shifter has been taken through
the differential because it is uniform. Hence
\begin{equation}
   \bG=\Grad\bw ,
   \tag{12.5}\label{eq:12.5}
\end{equation}
and $\bu=\shift\bw$. In components $w_{A}=\dd_{iA}u_{i}$,
$u_{i}=\dd_{iA}w_{A}$ and $G_{AB}=w_{A,B}$. Substituting
$\bH=\shift\bG$ into \eqr{7.22},
\begin{equation}
   \bF=\shift\big(\hI+\bG\big).
   \tag{12.6}\label{eq:12.6}
\end{equation}
The norms agree,
\[
   |\bG|^{2}=\tr\big(\bG\bG^{t}\big)
   =\tr\big(\shift^{t}\bH\bH^{t}\shift\big)
   =\tr\big(\bH\bH^{t}\shift\shift^{t}\big)=\tr\big(\bH\bH^{t}\big)
   =|\bH|^{2},
\]
by the cyclic property of the trace, so \eqr{12.1} is equivalent to
$|\bG|\ll1$, and again $|G_{AB}|\ll1$ componentwise.

Since $\bC$ is close to $\hI$, it is more convenient to work with the
Lagrange strain \eqr{2.97}, which measures departure from $\hI$
directly. From \eqr{12.3},
\begin{equation}
   \bE=\tfrac12\big(\bC-\hI\big)
   =\tfrac12\big(\bG+\bG^{t}+\bG^{t}\bG\big).
   \tag{12.7}\label{eq:12.7}
\end{equation}
Discarding the quadratic term,
\begin{equation}
   \bE\simeq\vek{\epsilon},\qquad\text{where}\quad
   \vek{\epsilon}=\tfrac12\big(\bG+\bG^{t}\big)
   \tag{12.8}\label{eq:12.8}
\end{equation}
is the \emph{referential infinitesimal strain tensor}, the symbol
$\simeq$ marking a statement valid to leading order in $|\bG|$.

\begin{remark}[small strain is weaker than small displacement
gradient]\label{rem:ch12weaker}
Equation \eqr{12.8} runs one way only. If $|\bG|\ll1$ then $|\bE|\ll1$,
since $|\bE|\le|\bG|+\tfrac12|\bG|^{2}$. The converse fails, and the
standard counterexample is a rigid rotation: take $\bF=\bQ$ with $\bQ$ a
finite rotation. Then $\bC=\hI$ and $\bE=\hOb$ exactly, so the strain is
not merely small but zero, while $\bG=\shift^{t}\bQ-\hI$ has norm of
order the rotation angle and is not small at all.

The whole theory therefore assumes more than small strain. It assumes
small \emph{rotation} as well, and \eqr{12.20}$_{3}$ makes that explicit
by producing a rotation tensor differing from the shifter by a term of
order $|\bG|$. A structure that strains by a thousandth while rotating
through a radian is outside the theory, which is why the analysis of
slender beams and thin shells needs more than what follows.
\end{remark}

\subsection*{Two observers, and why the shifter cannot be dropped}

These statements belong to the frame of one observer $O$. From
\eqr{12.3}$_{2}$ together with \eqr{7.18}, $\shift^{+}=\bQ\shift\bK^{t}$,
and \eqr{7.25}, $\bH^{+}=\bQ\bH\bK^{t}$,
\begin{equation}
   \bG^{+}=\bK\bG\bK^{t},
   \tag{12.9}\label{eq:12.9}
\end{equation}
$\bK\colon\hVt\to\hat{\mathbb{E}}^{3+}$ being the fixed orthogonal tensor relating the
two reference configurations. The computation is
\[
   \bG^{+}=(\shift^{+})^{t}\bH^{+}
   =\big(\bQ\shift\bK^{t}\big)^{t}\bQ\bH\bK^{t}
   =\bK\shift^{t}\bQ^{t}\bQ\bH\bK^{t}
   =\bK\shift^{t}\bH\bK^{t}=\bK\bG\bK^{t},
\]
$\bQ^{t}\bQ=\I$ removing the observer rotation entirely. Symmetrizing
\eqr{12.7} and \eqr{12.8} under conjugation, which commutes with
transposition because $(\bK\bA\bK^{t})^{t}=\bK\bA^{t}\bK^{t}$,
\begin{equation}
   \bE^{+}=\bK\bE\bK^{t},
   \tag{12.10}\label{eq:12.10}
\end{equation}
a result that also follows from \eqr{7.20} and \eqr{7.26} through
$\bE^{+}=\tfrac12[(\bF^{+})^{t}\bF^{+}-\hI^{+}]$. Further \eqr{12.9}
gives $|\bG^{+}|=|\bG|$, by the same cyclic-trace computation as above
with $\bK$ in place of $\shift$, and
\begin{equation}
   \vek{\epsilon}^{+}=\tfrac12\big[\bG^{+}+(\bG^{+})^{t}\big]
   =\bK\vek{\epsilon}\bK^{t},
   \tag{12.11}\label{eq:12.11}
\end{equation}
so that
\begin{equation}
   \bE^{+}\simeq\vek{\epsilon}^{+}.
   \tag{12.12}\label{eq:12.12}
\end{equation}
If one observer finds the strain small, so does every other, and
\eqr{12.12} says the approximation itself survives the change of
observer. This is Problem 1.

\begin{remark}[the equation that is not a tensor
equation]\label{rem:ch12shifter}
The literature almost always writes the strain--displacement relation by
dropping the shifter from \eqr{12.2}, obtaining
\[
   2\bE\simeq\bH+\bH^{t}
\]
in the frame of $O$. The left side is a referential tensor, mapping
$\hVt$ to $\hVt$. On the right, $\bH$ maps $\hVt$ to $\Vt$ and $\bH^{t}$
maps $\Vt$ to $\hVt$, so their sum is not defined: the two terms do not
even have the same domain, let alone the same codomain
(Figure~\ref{fig:ch12shifter}). The expression is not a tensor equation
and cannot be one.

The damage is visible under a change of observer. By \eqr{7.25},
\[
   \bH^{+}+(\bH^{+})^{t}=\bQ\bH\bK^{t}+\bK\bH^{t}\bQ^{t},
\]
which carries the observer rotation $\bQ$ on the outside, whereas
\eqr{12.10} requires $2\bE^{+}=\bK(2\bE)\bK^{t}$, with $\bK$ on both
sides and no $\bQ$ anywhere. So $\bH^{+}+(\bH^{+})^{t}$ cannot
approximate $2\bE^{+}$ for the second observer even if
$\bH+\bH^{t}$ did for the first. What makes \eqr{12.8} work is exactly
that $\bG=\shift^{t}\bH$ is referential on both legs, and \eqr{12.9}
then carries no $\bQ$.

None of this is an obstacle in practice, as the last paragraph of the
chapter records: an observer may choose the shifted basis
$\bE_{A}=\shift\hbE_{A}$ of \eqr{2.68} to coincide with $\be_{i}$, after
which $\dd_{iA}$ are numerically the Kronecker components and the
distinction is invisible. The point is that the invisibility is a
consequence of a choice, not a licence to write an equation that is
false in general.
\end{remark}

\begin{figure}[htbp]\centering
\renewcommand{\thefigure}{A}
\begin{tikzpicture}[scale=1.0,
   ar/.style={-{Stealth[length=2.4mm]},thick},
   nd/.style={inner sep=2pt}]
% ---- left: the illegal sum -----------------------------------------
\begin{scope}[xshift=0cm]
  \node[nd] (A) at (0,1.70)   {$\hVt$};
  \node[nd] (B) at (2.35,1.70) {$\Vt$};
  \draw[ar] (A) to[bend left=24]
     node[midway,above,font=\small] {$\bH$} (B);
  \draw[ar] (B) to[bend left=24]
     node[midway,below,font=\small] {$\bH^{t}$} (A);
  \node[font=\small,anchor=north,align=center] at (1.18,0.42)
     {$\bH+\bH^{t}$\\[1pt] has no domain};
\end{scope}
% ---- divider -------------------------------------------------------
\draw[gray!40] (4.05,-0.95)--(4.05,2.30);
% ---- right: the legal one ------------------------------------------
\begin{scope}[xshift=5.55cm]
  \node[nd] (C) at (0,1.70)   {$\hVt$};
  \node[nd] (D) at (2.35,1.70) {$\Vt$};
  \node[nd] (E) at (1.18,0.05) {$\hVt$};
  \draw[ar] (C)--(D) node[midway,above,font=\small] {$\bH$};
  \draw[ar] (D)--(E) node[midway,right=2pt,font=\small] {$\shift^{t}$};
  \draw[ar] (C)--(E) node[midway,left=2pt,font=\small] {$\bG$};
  \node[font=\small,anchor=north,align=center] at (1.18,-0.52)
     {$\bG+\bG^{t}$\\[1pt] maps $\hVt$ to $\hVt$};
\end{scope}
\end{tikzpicture}
\caption{Why Remark~\ref{rem:ch12shifter} insists on the shifter. The
displacement gradient $\bH$ of \eqr{7.22} is a two-point tensor, so its
transpose runs the other way and the sum $\bH+\bH^{t}$ is not defined
on any space. Composing with $\shift^{t}$ brings the second leg home,
and the resulting $\bG=\shift^{t}\bH$ is referential, so
$\bG+\bG^{t}$ is a legitimate referential tensor and is what
\eqr{12.8} uses. The triangle on the right commutes by the definition
\eqr{12.3}$_{2}$ of $\bG$.}
\label{fig:ch12shifter}
\end{figure}

\subsection*{The spatial description}

The displacement $\bu$ is a spatial vector, by \eqr{2.77} and
\eqr{2.78}, so it has a spatial gradient, defined by
$\dl\bu=(\grad\bu)\dl\bx$ at fixed $t$. Since also $\bH\dl\bX=\dl\bu$
and $\dl\bx=\bF\dl\bX$,
\begin{equation}
   \bH=(\grad\bu)\bF\simeq(\grad\bu)\shift ,
   \tag{12.13}\label{eq:12.13}
\end{equation}
the estimate by \eqr{12.6} and $|\bG|\ll1$. In components,
$\hat u_{i,A}=\tilde u_{i,j}F_{jA}=\tilde u_{i,j}\dd_{jB}
(\dd_{BA}+G_{BA})\simeq\tilde u_{i,j}\dd_{jA}$. This is an instance of
the general estimate
\begin{equation}
   (\cdot)_{,A}\simeq(\cdot)_{,j}\,\dd_{jA},
   \tag{12.14}\label{eq:12.14}
\end{equation}
valid to leading order: differentiating with respect to $\bX$ and with
respect to $\bx$ differ only by $\bF$, which is $\shift$ to leading
order. Combining \eqr{12.13} with \eqr{12.3}$_{2}$,
\begin{equation}
   \bG\simeq\shift^{t}(\grad\bu)\shift
   \qquad\text{and}\qquad
   \grad\bu\simeq\shift\bG\shift^{t},
   \tag{12.15}\label{eq:12.15}
\end{equation}
the second obtained from the first by multiplying by $\shift$ on the
left and $\shift^{t}$ on the right and using $\shift\shift^{t}=\I$.
Hence
\begin{equation}
   |\bH|=|\bG|\simeq|\grad\bu| ,
   \tag{12.16}\label{eq:12.16}
\end{equation}
so $|\tilde u_{i,j}|\ll1$ as well. Symmetrizing \eqr{12.15},
\begin{equation}
   \vek{\epsilon}\simeq\shift^{t}\vek{\varepsilon}\shift,
   \qquad\text{where}\quad
   \vek{\varepsilon}=\tfrac12\big[\grad\bu+(\grad\bu)^{t}\big]
   \tag{12.17}\label{eq:12.17}
\end{equation}
is the \emph{spatial infinitesimal strain tensor}, with components
satisfying $\epsilon_{AB}\simeq\dd_{iA}\dd_{jB}\varepsilon_{ij}$. From
the exact relation \eqr{12.13}$_{1}$, or from either estimate,
\begin{equation}
   \grad^{+}\bu^{+}=\bQ(\grad\bu)\bQ^{t}
   \qquad\text{and}\qquad
   \vek{\varepsilon}^{+}=\bQ\vek{\varepsilon}\bQ^{t},
   \tag{12.18}\label{eq:12.18}
\end{equation}
$\bQ\colon\Vt\to\mathbb{E}^{3+}$ being the observer rotation. These also follow
from \eqr{7.3} and \eqr{7.21}. Note the contrast with \eqr{12.9} and
\eqr{12.11}: the spatial objects are conjugated by $\bQ$, the
referential ones by $\bK$, and the shifter is what carries one rule into
the other.

\subsection*{Stretch and rotation}

Take the polar decomposition $\bF=\bR\bU$ of \eqr{3.32}. From
\eqr{12.3} and \eqr{12.8},
\begin{equation}
   \bU^{2}-\hI=\bC-\hI=2\bE\simeq2\vek{\epsilon} .
   \tag{12.19}\label{eq:12.19}
\end{equation}
Write $\bU=\hI+\vek{\beta}$, which defines $\vek{\beta}$; then
$\bU^{2}=\hI+2\vek{\beta}+\vek{\beta}^{2}$, and discarding the square,
$2\vek{\beta}\simeq2\vek{\epsilon}$. Hence, with
$\bU^{-1}=(\hI+\vek{\beta})^{-1}\simeq\hI-\vek{\beta}$,
\begin{equation}
   \bU-\hI\simeq\vek{\epsilon},\qquad
   \bU^{-1}-\hI\simeq-\vek{\epsilon},\qquad
   \bR-\shift\simeq\shift\hat{\vek{\omega}},
   \tag{12.20}\label{eq:12.20}
\end{equation}
where
\begin{equation}
   \hat{\vek{\omega}}=\tfrac12\big(\bG-\bG^{t}\big)
   \tag{12.21}\label{eq:12.21}
\end{equation}
is the \emph{referential infinitesimal rotation tensor}. The third of
\eqr{12.20} is the one that needs a computation. From $\bR=\bF\bU^{-1}$
with \eqr{12.6} and $\bU^{-1}\simeq\hI-\vek{\epsilon}$,
\[
   \bR\simeq\shift\big(\hI+\bG\big)\big(\hI-\vek{\epsilon}\big)
   \simeq\shift\big(\hI+\bG-\vek{\epsilon}\big)
   =\shift\big(\hI+\hat{\vek{\omega}}\big),
\]
the product $\bG\vek{\epsilon}$ being of second order, and
$\bG-\vek{\epsilon}=\bG-\tfrac12(\bG+\bG^{t})
=\tfrac12(\bG-\bG^{t})=\hat{\vek{\omega}}$. Equivalently,
\begin{equation}
   \bR-\shift\simeq\vek{\omega}\shift,\qquad\text{where}\quad
   \vek{\omega}=\tfrac12\big[\grad\bu-(\grad\bu)^{t}\big]
   =\shift\hat{\vek{\omega}}\shift^{t}
   \tag{12.22}\label{eq:12.22}
\end{equation}
is the \emph{spatial infinitesimal rotation tensor}, the two forms
agreeing because
$\vek{\omega}\shift=\shift\hat{\vek{\omega}}\shift^{t}\shift
=\shift\hat{\vek{\omega}}$.

\begin{remark}[the linearized polar decomposition]%
\label{rem:ch12polar}
Collecting \eqr{12.8}, \eqr{12.21} and \eqr{12.20},
\[
   \bG=\vek{\epsilon}+\hat{\vek{\omega}},\qquad
   \bF=\shift\big(\hI+\vek{\epsilon}+\hat{\vek{\omega}}\big)
   \simeq\shift\big(\hI+\hat{\vek{\omega}}\big)
   \big(\hI+\vek{\epsilon}\big),
\]
the last step discarding
$\hat{\vek{\omega}}\vek{\epsilon}$. So at leading order the polar
decomposition degenerates into the additive split of $\bG$ into its
symmetric and skew parts, and the multiplicative structure that
Chapter 3 worked so hard for collapses into a sum. That collapse is the
whole simplification of linear elasticity, and
Remark~\ref{rem:ch12weaker} is the price: the additive split is
faithful only while the rotation is small, since
$\hI+\hat{\vek{\omega}}$ is a rotation only to first order.

The additive split is also the linear analogue of \eqr{4.7}. There
$\grad\bv$ split into stretching and spin; here $\grad\bu$ splits into
strain and infinitesimal rotation. Chapter 4's split is exact and
Chapter 12's is not, because velocity is a rate and displacement is
not.
\end{remark}

\subsection*{Volume and density}

The local volume change is $J-1$, and $J^{2}=\det\bC$. Writing
$\bC=\hI+\bA$ with $\bA=\bG+\bG^{t}+\bG^{t}\bG$ from \eqr{12.3}, and
using the result of Problem 23 of Chapter 1 for the expansion of a
determinant,
\begin{equation}
   J^{2}=\det\big(\hI+\bA\big)
   =\det\big(\bA-\lambda\hI\big)\big|_{\lambda=-1},
   \tag{12.23}\label{eq:12.23}
\end{equation}
whose expansion is $1+I_{1}(\bA)+I_{2}(\bA)+I_{3}(\bA)$. The last two
invariants are quadratic and cubic in $\bA$ and so are discarded, and
$\tr(\bG^{t}\bG)$ is quadratic in $\bG$, leaving
\begin{equation}
   J^{2}-1\simeq\tr\bA\simeq2\tr\bG\simeq2\tr(\grad\bu),
   \tag{12.24}\label{eq:12.24}
\end{equation}
the last estimate by \eqr{12.15} and
$\tr[\shift^{t}(\grad\bu)\shift]=\tr[(\grad\bu)\shift\shift^{t}]
=\tr[(\grad\bu)\I]=\tr(\grad\bu)$. Since $J$ is close to $1$,
$J^{2}-1=(J-1)(J+1)\simeq2(J-1)$, and therefore
\begin{equation}
   J-1\simeq\Dvg\bw\simeq\dvg\bu .
   \tag{12.25}\label{eq:12.25}
\end{equation}

The change of volume of an arbitrary part $S$ follows by integrating,
$P$ being the region occupied by $S$ in $\kappa$:
\begin{equation}
   \Delta V=\int_{P}\big(J-1\big)\,\dl V\simeq\int_{P}\Dvg\bw\,\dl V
   =\int_{\partial P}\ip\bw\bN\,\dl A ,
   \tag{12.26}\label{eq:12.26}
\end{equation}
by \eqr{2.50} and the divergence theorem. Alternatively, from
$J^{-1}-1\simeq-\dvg\bu$,
\begin{equation}
   \Delta V\simeq\int_{P}\dvg\bu\,\dl V
   =\int_{P_{t}}J^{-1}\dvg\bu\,\dl v
   \simeq\int_{P_{t}}\dvg\bu\,\dl v=\int_{\partial P_{t}}\ip\bu\bn\,\dl a .
   \tag{12.27}\label{eq:12.27}
\end{equation}
To leading order the volume change is the flux of the displacement
through the boundary, computed with either description. That the two
right-hand sides agree is Problem 4.

Finally the density. By \eqr{2.59}, $\rho_{\kappa}=\rho J$, so
\begin{equation}
   \Delta\rho=\rho-\rho_{\kappa},
   \tag{12.28}\label{eq:12.28}
\end{equation}
and $\rho=\rho_{\kappa}J^{-1}\simeq\rho_{\kappa}(1-\Dvg\bw)$ gives
\begin{equation}
   \Delta\rho\simeq-\rho_{\kappa}\Dvg\bw\simeq-\rho_{\kappa}\dvg\bu,
   \qquad\text{yielding}\qquad \rho\simeq\rho_{\kappa}
   \tag{12.29}\label{eq:12.29}
\end{equation}
at leading order. The density change is of first order in $|\bG|$ and
the density itself is $\rho_{\kappa}$ to zeroth order, which is why the
linear equation of motion may carry either symbol.

%=======================================================================
\section{Stress--deformation relation}\label{sec:ch12stress}
%=======================================================================

We now want small-displacement-gradient estimates of the three stresses.
Introduce the \emph{strain-energy function}
\begin{equation}
   U(\bE;\bX)=\bar W\big(\theta_{0},\hI+2\bE;\bX\big),
   \tag{12.30}\label{eq:12.30}
\end{equation}
$\bar W$ being the free energy per unit reference volume of \eqr{11.19}
and $\theta_{0}$ the fixed uniform temperature. The substitution
$\bC=\hI+2\bE$ is \eqr{2.97} read backwards. Proceeding as at
\eqr{11.20}, and using $\dot\bC=2\dot\bE$,
\begin{equation}
   U_{\bE}\cdot\dot\bE=\bar W_{\bC}\cdot\dot\bC
   =2\bar W_{\bC}\cdot\dot\bE=\bS\cdot\dot\bE
   \tag{12.31}\label{eq:12.31}
\end{equation}
for every symmetric $\dot\bE$, the last step by \eqr{11.24}. Since
$U_{\bE}$ is symmetric, being a gradient with respect to a symmetric
argument, and $\dot\bE$ ranges over all symmetric tensors,
\begin{equation}
   \bS=U_{\bE} .
   \tag{12.32}\label{eq:12.32}
\end{equation}

\subsection*{Expanding the energy}

A small displacement gradient implies a small strain, by
Remark~\ref{rem:ch12weaker}, so we want $U$ near $\bE=\hOb$. Taylor's
theorem gives
\begin{equation}
   U(\bE;\bX)=U(\hOb;\bX)+\bS_{0}(\bX)\cdot\bE
   +\tfrac12\bE\cdot\mathbb{D}(\bX)[\bE]+o\big(|\bE|^{2}\big),
   \tag{12.33}\label{eq:12.33}
\end{equation}
where
\begin{equation}
   \bS_{0}(\bX)=U_{\bE}\big|_{\bE=\hOb}
   \qquad\text{and}\qquad
   \mathbb{D}(\bX)=U_{\bE\bE}\big|_{\bE=\hOb}
   \tag{12.34}\label{eq:12.34}
\end{equation}
is a referential fourth-order tensor, $\mathbb{D}(\bX)[\bE]$ being a
linear second-order-tensor-valued function of $\bE$.

\begin{remark}[$\mathbb{C}$ here is not the complex
numbers]\label{rem:ch12Cclash}
The spatial counterpart of $\mathbb{D}$, introduced at \eqr{12.56}, is
written $\mathbb{C}$, following the book. Chapter 1 used the same
blackboard letter for the complex numbers in the proof that the
principal values of a symmetric tensor are real. The two never meet:
this chapter has no complex numbers in it. Nor should $\mathbb{C}$ be
confused with $\bC$, the right Cauchy--Green tensor, which carries an
under-tilde and is of second order.
\end{remark}

Resolving on a basis, $\bE=E_{CD}\hbE_{C}\otimes\hbE_{D}$, linearity
gives
\begin{equation}
   \mathbb{D}(\bX)[\bE]
   =E_{CD}\,\mathbb{D}(\bX)\big[\hbE_{C}\otimes\hbE_{D}\big],
   \tag{12.35}\label{eq:12.35}
\end{equation}
and therefore
\begin{equation}
\begin{aligned}
   \bE\cdot\mathbb{D}(\bX)[\bE]&=E_{AB}E_{CD}\,\mathbb{D}_{ABCD}(\bX),\\
   \text{where}\quad \mathbb{D}_{ABCD}
   &=\hbE_{A}\otimes\hbE_{B}\cdot
     \mathbb{D}\big[\hbE_{C}\otimes\hbE_{D}\big]
   =\hbE_{A}\cdot\big(\mathbb{D}\big[\hbE_{C}\otimes\hbE_{D}\big]\big)
     \hbE_{B}.
\end{aligned}
   \tag{12.36}\label{eq:12.36}
\end{equation}

Three symmetries may now be assumed without loss of generality, because
only the combination $E_{AB}E_{CD}\mathbb{D}_{ABCD}$ ever appears.
Relabelling the summation indices, and using that ordinary
multiplication commutes,
\begin{equation}
   E_{AB}E_{CD}\mathbb{D}_{ABCD}=E_{CD}E_{AB}\mathbb{D}_{CDAB}
   =E_{AB}E_{CD}\mathbb{D}_{CDAB},
   \tag{12.37}\label{eq:12.37}
\end{equation}
so nothing is lost by taking
\begin{equation}
   \mathbb{D}_{ABCD}=\mathbb{D}_{CDAB} .
   \tag{12.38}\label{eq:12.38}
\end{equation}
Relabelling again and using $E_{AB}=E_{BA}$,
\begin{equation}
   E_{BA}E_{CD}\mathbb{D}_{BACD}=E_{AB}E_{CD}\mathbb{D}_{ABCD}
   =E_{BA}E_{CD}\mathbb{D}_{ABCD},
   \tag{12.39}\label{eq:12.39}
\end{equation}
so we may take
\begin{equation}
   \mathbb{D}_{BACD}=\mathbb{D}_{ABCD}
   \tag{12.40}\label{eq:12.40}
\end{equation}
and hence, combining with \eqr{12.38},
\begin{equation}
   \mathbb{D}_{ABCD}=\mathbb{D}_{ABDC} .
   \tag{12.41}\label{eq:12.41}
\end{equation}

In tensor form, let $A_{AB}$ and $B_{CD}$ be the components of arbitrary
referential tensors. Then \eqr{12.36}$_{2}$, \eqr{12.38} and linearity
give the \emph{major symmetry}
\begin{equation}
   \bA\cdot\mathbb{D}[\bB]=A_{AB}B_{CD}\mathbb{D}_{ABCD}
   =B_{CD}A_{AB}\mathbb{D}_{CDAB}=\bB\cdot\mathbb{D}[\bA]
   \tag{12.42}\label{eq:12.42}
\end{equation}
for all $\bA$ and $\bB$, while \eqr{12.40} and \eqr{12.41} give the two
\emph{minor symmetries}
\begin{equation}
   \bA^{t}\cdot\mathbb{D}[\bB]=\bA\cdot\mathbb{D}[\bB]
   \qquad\text{and}\qquad
   \bA\cdot\mathbb{D}\big[\bB^{t}\big]=\bA\cdot\mathbb{D}[\bB].
   \tag{12.43}\label{eq:12.43}
\end{equation}

Now differentiate \eqr{12.33}, keeping terms through quadratic order.
The rate of the linear map is
\begin{equation}
   \big(\mathbb{D}(\bX)[\bE]\big)^{\textstyle\cdot}
   =\dot E_{CD}\,\mathbb{D}(\bX)\big[\hbE_{C}\otimes\hbE_{D}\big]
   =\mathbb{D}(\bX)\big[\dot\bE\big],
   \tag{12.44}\label{eq:12.44}
\end{equation}
by \eqr{12.35}, the basis dyads being fixed. Hence, using major
symmetry to merge the two quadratic contributions,
\begin{equation}
\begin{aligned}
   U_{\bE}\cdot\dot\bE
   &=\bS_{0}(\bX)\cdot\dot\bE
   +\tfrac12\big\{\bE\cdot\mathbb{D}(\bX)[\dot\bE]
   +\dot\bE\cdot\mathbb{D}(\bX)[\bE]\big\}\\
   &=\big\{\bS_{0}(\bX)+\mathbb{D}(\bX)[\bE]\big\}\cdot\dot\bE ,
\end{aligned}
   \tag{12.45}\label{eq:12.45}
\end{equation}
for all symmetric $\dot\bE$, and \eqr{12.32} with \eqr{12.40} delivers
the Piola--Kirchhoff stress
\begin{equation}
   \bS=\bS_{0}(\bX)+\mathbb{D}(\bX)[\bE].
   \tag{12.46}\label{eq:12.46}
\end{equation}
Here $\bS_{0}$ is the \emph{residual stress}, the stress carried at zero
strain, and $\mathbb{D}$ is the \emph{stiffness tensor}, also called the
elastic modulus tensor.

\subsection*{Isotropy}

For an isotropic material, \eqr{11.29} gives
\begin{equation}
\begin{aligned}
   U(\bE;\bX)&=\bar W\big(\theta_{0},\hI+2\bE;\bX\big)
   =\bar W\big(\theta_{0},\bR^{t}(\hI+2\bE)\bR;\bX\big)\\
   &=\bar W\big(\theta_{0},\hI+2\bR^{t}\bE\bR;\bX\big)
   =U\big(\bR^{t}\bE\bR;\bX\big)
\end{aligned}
   \tag{12.47}\label{eq:12.47}
\end{equation}
for all $\bR\in\Orth$, the third step by $\bR^{t}\hI\bR=\bR^{t}\bR=\hI$.
So $U$ is an isotropic scalar function of $\bE$ and depends on it
through the principal invariants $I_{1}(\bE)=\tr\bE=\hI\cdot\bE$,
$I_{2}(\bE)=\tfrac12[I^{2}_{1}-\tr(\bE^{2})]$ and $I_{3}(\bE)=\det\bE$,
or equivalently through $\hI\cdot\bE$, $\bE\cdot\bE$ and $\det\bE$.
Keeping terms through quadratic order, and noting that $\det\bE$ is
cubic while $I_{2}$ is quadratic,
\begin{equation}
   U(\bE;\bX)=U(\hOb;\bX)+\alpha(\bX)\,\hI\cdot\bE
   +\tfrac12\lambda(\bX)\big(\hI\cdot\bE\big)^{2}
   +\mu(\bX)\,\bE\cdot\bE ,
   \tag{12.48}\label{eq:12.48}
\end{equation}
the functions $\alpha,\lambda,\mu$ being properties of the material.
Differentiating term by term,
\begin{equation}
   U_{\bE}\cdot\dot\bE
   =\big\{\big[\alpha+\lambda(\tr\bE)\big]\hI+2\mu\bE\big\}\cdot\dot\bE ,
   \tag{12.49}\label{eq:12.49}
\end{equation}
since $(\hI\cdot\bE)^{\textstyle\cdot}=\hI\cdot\dot\bE$ and
$(\bE\cdot\bE)^{\textstyle\cdot}=2\bE\cdot\dot\bE$, and comparison with
\eqr{12.46} gives
\begin{equation}
   \bS_{0}(\bX)=\alpha(\bX)\hI
   \qquad\text{and}\qquad
   \mathbb{D}(\bX)[\bE]=\lambda(\bX)(\tr\bE)\hI+2\mu(\bX)\bE .
   \tag{12.50}\label{eq:12.50}
\end{equation}

\begin{remark}[what isotropy costs the residual
stress]\label{rem:ch12resid}
Equation \eqr{12.50}$_{1}$ says that an isotropic material can carry
only a hydrostatic residual stress. That is the linearized form of the
statement proved at Problem 10 of Chapter 11, and Problem 3 below turns
it into the same conclusion: with equilibrium and one traction-free
patch, an isotropic body cannot be residually stressed at all. The two
constants $\lambda$ and $\mu$ are the Lam\'e moduli, and $\mu$ is the
shear modulus $G$ of \eqr{11.41}.
\end{remark}

\subsection*{Linearizing, and the other two stresses}

Equation \eqr{12.46} is accurate whenever $|\bE|$ is small, and it is a
nonlinear relation between $\bS$ and $\bG$, since $\bE$ is quadratic in
$\bG$ by \eqr{12.7}. If $|\bG|$ is also small, \eqr{12.8} linearizes it:
\begin{equation}
   \bS-\bS_{0}(\bX)\simeq\mathbb{D}(\bX)[\vek{\epsilon}]
   =\mathbb{D}(\bX)[\bG],
   \tag{12.51}\label{eq:12.51}
\end{equation}
the second equality by the minor symmetry \eqr{12.43}$_{2}$, which lets
the skew part of $\bG$ be discarded; in components
$S_{AB}-S_{(0)AB}\simeq\mathbb{D}_{ABCD}w_{C,D}$. Classical linear
elasticity assumes the residual stress vanishes, so
\begin{equation}
   \bS\simeq\mathbb{D}(\bX)[\bG],
   \tag{12.52}\label{eq:12.52}
\end{equation}
and for an isotropic material, by \eqr{12.50}$_{2}$ with
$\vek{\epsilon}=\tfrac12(\bG+\bG^{t})$,
\begin{equation}
   \bS\simeq\lambda(\bX)(\tr\bG)\hI+\mu(\bX)\big(\bG+\bG^{t}\big).
   \tag{12.53}\label{eq:12.53}
\end{equation}

For the Piola stress, combine \eqr{6.144} with \eqr{12.6}, which gives
$\bF\simeq\shift$:
\begin{equation}
   \bP=\bF\bS\simeq\shift\bS .
   \tag{12.54}\label{eq:12.54}
\end{equation}
To the same order $\bF^{*}\simeq\shift$, which is Problem 4, and either
\eqr{6.140} or \eqr{6.145} then gives
\begin{equation}
   \bT\simeq\shift\bS\shift^{t}
   \tag{12.55}\label{eq:12.55}
\end{equation}
for the Cauchy stress. In components $P_{iA}\simeq\dd_{iB}S_{BA}$ and
$T_{ij}\simeq\dd_{iA}\dd_{jB}S_{AB}$: all three stresses carry the same
numbers, and differ only in which legs are referential.

With vanishing residual stress, \eqr{12.15} and \eqr{12.52} combine into
\begin{equation}
   \bT\simeq\mathbb{C}[\grad\bu],\qquad\text{where}\quad
   \mathbb{C}[\grad\bu]
   =\shift\big\{\mathbb{D}\big[\shift^{t}(\grad\bu)\shift\big]\big\}
   \shift^{t}.
   \tag{12.56}\label{eq:12.56}
\end{equation}
The definition is easier to read in components. Using
\eqr{12.36}$_{2}$ and linearity at each step,
\begin{equation}
\begin{aligned}
   T_{ij}&=\ip{\be_{i}}{\bT\be_{j}}
   =\ip{\shift^{t}\be_{i}}
     {\big(\mathbb{D}\big[u_{k,l}\dd_{kC}\dd_{lD}
      \hbE_{C}\otimes\hbE_{D}\big]\big)\shift^{t}\be_{j}}\\
   &=\dd_{iA}\dd_{jB}\,\hbE_{A}\cdot
     \big(\mathbb{D}\big[\dd_{kC}\dd_{lD}u_{k,l}
     \hbE_{C}\otimes\hbE_{D}\big]\big)\hbE_{B}\\
   &=\dd_{iA}\dd_{jB}\dd_{kC}\dd_{lD}\,u_{k,l}\,
     \hbE_{A}\cdot\big(\mathbb{D}\big[\hbE_{C}\otimes\hbE_{D}\big]\big)
     \hbE_{B}
   =\mathbb{C}_{ijkl}\,u_{k,l},
\end{aligned}
   \tag{12.57}\label{eq:12.57}
\end{equation}
where
\begin{equation}
   \mathbb{C}_{ijkl}
   =\dd_{iA}\dd_{jB}\dd_{kC}\dd_{lD}\,\mathbb{D}_{ABCD}.
   \tag{12.58}\label{eq:12.58}
\end{equation}
The spatial stiffness inherits every symmetry of the referential one,
which is Problem 5: the major symmetry
\begin{equation}
   \mathbb{C}_{ijkl}=\mathbb{C}_{klij}
   \tag{12.59}\label{eq:12.59}
\end{equation}
and the two minor ones
\begin{equation}
   \mathbb{C}_{ijkl}=\mathbb{C}_{jikl}
   \qquad\text{and}\qquad
   \mathbb{C}_{ijkl}=\mathbb{C}_{ijlk}.
   \tag{12.60}\label{eq:12.60}
\end{equation}
The last of these lets $u_{k,l}$ in \eqr{12.57} be replaced by its
symmetric part, so $T_{ij}=\mathbb{C}_{ijkl}\varepsilon_{kl}$ with
$\varepsilon_{ij}$ the components of \eqr{12.17}. For an isotropic
material, \eqr{12.15} and \eqr{12.53} give directly
\begin{equation}
   \bT\simeq\lambda(\dvg\bu)\I
   +\mu\big[\grad\bu+(\grad\bu)^{t}\big].
   \tag{12.61}\label{eq:12.61}
\end{equation}

\subsection*{Tractions}

The Piola traction is $\bp=\bP\bN$ and the Cauchy traction $\bt=\bT\bn$,
by \eqr{6.132} and \eqr{6.94}. The Piola--Nanson formula \eqr{2.54} with
$\bF^{*}\simeq\shift$ gives
\begin{equation}
   \bn\,\dl a\simeq\shift\bN\,\dl A ,
   \tag{12.62}\label{eq:12.62}
\end{equation}
so $\dl a\simeq\dl A$, the shifter preserving length, and
$\bn\simeq\shift\bN$. Therefore
\begin{equation}
   \bp=\bP\bN=\bT\bF^{*}\bN\simeq\bT\shift\bN\simeq\bT\bn=\bt :
   \tag{12.63}\label{eq:12.63}
\end{equation}
the two tractions coincide at leading order, which is why elementary
treatments need not distinguish them. By \eqr{12.54} this is the same
as
\begin{equation}
   \shift^{t}\bp\simeq\bS\bN .
   \tag{12.64}\label{eq:12.64}
\end{equation}

In the frame of $O^{+}$ the counterparts follow from Chapter 7:
\begin{equation}
\begin{aligned}
   \bS^{+}&\simeq\mathbb{D}^{+}[\bG^{+}],\qquad
   \bT^{+}\simeq\shift^{+}\bS^{+}(\shift^{+})^{t}
   \simeq\mathbb{C}^{+}[\grad^{+}\bu^{+}],\\
   (\shift^{+})^{t}\bp^{+}&\simeq\bS^{+}\bN^{+},
\end{aligned}
   \tag{12.65}\label{eq:12.65}
\end{equation}
where
\begin{equation}
\begin{aligned}
   \mathbb{D}^{+}[\bG^{+}]&=\bK\big(\mathbb{D}[\bK^{t}\bG^{+}\bK]\big)
   \bK^{t},\\
   \mathbb{C}^{+}[\grad^{+}\bu^{+}]
   &=\bQ\big\{\mathbb{C}\big[\bQ^{t}(\grad^{+}\bu^{+})\bQ\big]\big\}
   \bQ^{t}.
\end{aligned}
   \tag{12.66}\label{eq:12.66}
\end{equation}
For an isotropic material these reduce to
\begin{equation}
\begin{aligned}
   \mathbb{D}^{+}[\bG^{+}]
   &=\lambda(\tr\bG^{+})\hI^{+}+\mu\big[\bG^{+}+(\bG^{+})^{t}\big],\\
   \mathbb{C}^{+}[\grad^{+}\bu^{+}]
   &=\lambda(\dvg^{+}\bu^{+})\I^{+}
   +\mu\big[\grad^{+}\bu^{+}+\big(\grad^{+}\bu^{+}\big)^{t}\big],
\end{aligned}
   \tag{12.67}\label{eq:12.67}
\end{equation}
with the same two constants: the moduli are numbers and no observer can
change them.

%=======================================================================
\section{Equation of motion}\label{sec:ch12eom}
%=======================================================================

Take the frame of $O$ to be inertial. The referential equation of motion
is \eqr{6.136} with $\bP=\bF\bS$ and, the shifter being fixed,
\begin{equation}
   \dot\bv=\ddot{\bchi}_{\kappa}
   =\big(\bu+\shift\bX\big)^{\textstyle\cdot\cdot}=\ddot\bu .
   \tag{12.68}\label{eq:12.68}
\end{equation}
The divergence of the Piola stress is $\Dvg\bP=P_{iA,A}\be_{i}$, and by
the product rule together with $F_{iB}=\dd_{iB}+\dd_{iC}G_{CB}$ from
\eqr{12.6},
\begin{equation}
\begin{aligned}
   P_{iA,A}&=\big(F_{iB}S_{BA}\big)_{,A}
   =F_{iB}S_{BA,A}+\dd_{iC}G_{CB,A}S_{BA}\\
   &\simeq\dd_{iB}S_{BA,A}+\dd_{iC}G_{CB,A}S_{BA}.
\end{aligned}
   \tag{12.69}\label{eq:12.69}
\end{equation}
With no residual stress, \eqr{12.52} makes the non-dimensionalized
$S_{BA}$ of order $\varepsilon$, the small number of \eqr{12.1}. If we
assume in addition that $G_{CB,A}$, non-dimensionalized by a length
scale of the problem, is also of order $\varepsilon$, the second term is
of order $\varepsilon^{2}$ and drops:
\begin{equation}
   \Dvg\bP\simeq\shift\big(\Dvg\bS\big).
   \tag{12.70}\label{eq:12.70}
\end{equation}

\begin{remark}[the assumption hidden in \eqr{12.70}]%
\label{rem:ch12extra}
Smallness of $\bG$ says nothing about smallness of $\Grad\bG$. A field
may be uniformly small and vary arbitrarily rapidly, and then the
discarded term in \eqr{12.69} is not negligible. The restriction on
$G_{CB,A}$ is therefore a second and independent hypothesis, tacitly
made throughout the literature. It is what confines the theory to
displacement fields whose variation is slow on the scale of the body,
and it excludes, for instance, the immediate neighbourhood of a crack
tip, where the strain is small but its gradient is not.

The same caution applies to the solution rather than the data. Once a
boundary-value problem has been solved, both hypotheses should be
checked against the answer, since nothing in the linear theory enforces
them.
\end{remark}

With these estimates, and applying $\shift^{t}$ to
$\shift\Dvg\bS+\rho_{\kappa}\bb=\rho_{\kappa}\ddot\bu$,
\begin{equation}
   \Dvg\bS+\rho_{\kappa}\bff=\rho_{\kappa}\ddot\bw ,
   \tag{12.71}\label{eq:12.71}
\end{equation}
where $\bff=\shift^{t}\bb$ is the referential body force,
$\ddot\bw=(\shift^{t}\bu)^{\textstyle\cdot\cdot}=\shift^{t}\ddot\bu$ the
referential acceleration, and
\begin{equation}
   \bS=\mathbb{D}[\Grad\bw].
   \tag{12.72}\label{eq:12.72}
\end{equation}
This is a linear second-order partial differential system for
$\bw(\bX,t)$. A typical initial-boundary-value problem adds initial data
$\bw|_{t=0}$ and $\dot\bw|_{t=0}$, together with $\shift^{t}\bp$
prescribed on a part $\partial\kappa_{p}$ of the boundary and $\bw$
prescribed on the complement
$\partial\kappa_{w}=\partial\kappa\setminus\partial\kappa_{p}$.

In the frame of $O^{+}$,
\begin{equation}
   \Dvg^{+}\bS^{+}=\bK\big(\Dvg\bS\big),
   \tag{12.73}\label{eq:12.73}
\end{equation}
which is Problem 9, and hence, with \eqr{7.38}$_{3}$,
\begin{equation}
   \Dvg^{+}\bS^{+}+\rho^{+}_{\kappa^{+}}\bff^{+}
   =\rho^{+}_{\kappa^{+}}\bK\ddot\bw ,
   \tag{12.74}\label{eq:12.74}
\end{equation}
with
\begin{equation}
   \bff^{+}=(\shift^{+})^{t}\bb^{+}
   \qquad\text{and}\qquad
   \bS^{+}=\mathbb{D}^{+}[\Grad^{+}\bw^{+}],\quad
   \bw^{+}=(\shift^{+})^{t}\bu^{+}.
   \tag{12.75}\label{eq:12.75}
\end{equation}
If the frame of $O^{+}$ is also inertial then the two are related by a
Galilean transformation \eqr{7.10}, with $\bQ$ a fixed orthogonal tensor
and $\bV$ a fixed vector, and then $\bK\ddot\bw=\ddot\bw^{+}$, so that
\begin{equation}
   \Dvg^{+}\bS^{+}+\rho^{+}_{\kappa^{+}}\bff^{+}
   =\rho^{+}_{\kappa^{+}}\ddot\bw^{+}.
   \tag{12.76}\label{eq:12.76}
\end{equation}
Linear elasticity therefore holds in every inertial frame, in the same
form. The non-inertial case is Problem 10.

\begin{remark}[how to make the shifter disappear]%
\label{rem:ch12disappear}
The shifter has done real work: without it, \eqr{12.8} is not a tensor
equation and \eqr{12.9} does not follow, as
Remark~\ref{rem:ch12shifter} showed. In any single problem, however, it
can be made invisible. The observer $O$ is free to choose the shifted
basis $\bE_{A}=\shift\hbE_{A}$ of \eqr{2.68} to coincide with the
spatial basis $\be_{i}$. Then $\dd_{iA}$ are numerically the components
of the Kronecker delta, the components of a referential tensor and of
its spatial shift are the same numbers, and one may write
$2\bE\simeq\bH+\bH^{t}$ without harm.

That is why the omission survives in the literature. It is a legitimate
computational convenience for one observer, and an illegitimate
statement about tensors. The distinction matters exactly when a second
observer appears, which is the situation \eqr{12.65} and \eqr{12.66}
describe.
\end{remark}

%=======================================================================
\setcounter{problemenv}{0}
\section{Problems}\label{sec:problems:c12}
%=======================================================================

\begin{problem}[The kinematic transformations]
Derive \eqr{12.9} to \eqr{12.11}. Show that $|\bG^{+}|=|\bG|$. What is
the relationship between $\bw^{+}$ and $\bw$?
\end{problem}

\begin{solution}
\textbf{\eqr{12.9}.} By definition $\bG^{+}=(\shift^{+})^{t}\bH^{+}$.
Chapter 7 supplies both factors: $\shift^{+}=\bQ\shift\bK^{t}$ by
\eqr{7.18} and $\bH^{+}=\bQ\bH\bK^{t}$ by \eqr{7.25}. Hence
\[
   \bG^{+}=\big(\bQ\shift\bK^{t}\big)^{t}\,\bQ\bH\bK^{t}
   =\bK\shift^{t}\bQ^{t}\bQ\bH\bK^{t}
   =\bK\shift^{t}\bH\bK^{t}
   =\boxed{\;\bK\bG\bK^{t}\;}
\]
using $(\bA\bB\bC)^{t}=\bC^{t}\bB^{t}\bA^{t}$ and $\bQ^{t}\bQ=\I$. The
observer rotation cancels completely, which is the point: $\bG$ is a
referential tensor and only $\bK$ can act on it.

\textbf{\eqr{12.10}.} Substituting into \eqr{12.7} and using
$(\bK\bG\bK^{t})^{t}=\bK\bG^{t}\bK^{t}$,
\[
\begin{aligned}
   \bE^{+}&=\tfrac12\big[\bG^{+}+(\bG^{+})^{t}
   +(\bG^{+})^{t}\bG^{+}\big]\\
   &=\tfrac12\big[\bK\bG\bK^{t}+\bK\bG^{t}\bK^{t}
   +\bK\bG^{t}\bK^{t}\bK\bG\bK^{t}\big]
   =\bK\bE\bK^{t},
\end{aligned}
\]
the last step by $\bK^{t}\bK=\hI$ in the quadratic term and by
factoring $\bK$ out on the left and $\bK^{t}$ on the right. So
\eqr{12.10} is exact, not an approximation.

\textbf{\eqr{12.11}.} The same computation without the quadratic term
gives $\vek{\epsilon}^{+}=\tfrac12[\bG^{+}+(\bG^{+})^{t}]
=\bK\vek{\epsilon}\bK^{t}$. Comparing with \eqr{12.10}, the two agree
to leading order exactly when $\bE\simeq\vek{\epsilon}$ does, which is
\eqr{12.12}: the approximation is preserved under change of observer.

\textbf{The norm.} Using the cyclic property of the trace and
$\bK^{t}\bK=\hI$ twice,
\[
\begin{aligned}
   |\bG^{+}|^{2}&=\tr\big(\bG^{+}(\bG^{+})^{t}\big)
   =\tr\big(\bK\bG\bK^{t}\bK\bG^{t}\bK^{t}\big)\\
   &=\tr\big(\bG\bG^{t}\bK^{t}\bK\big)
   =\tr\big(\bG\bG^{t}\big)=|\bG|^{2},
\end{aligned}
\]
so $\boxed{|\bG^{+}|=|\bG|}$. The two observers agree on how small the
displacement gradient is, and therefore on whether the linear theory
applies at all.

\textbf{The displacements.} By \eqr{7.21}, $\bu^{+}=\bQ\bu+\bd(t)$ with
$\bd$ spatially uniform. Hence
\[
   \bw^{+}=(\shift^{+})^{t}\bu^{+}
   =\bK\shift^{t}\bQ^{t}\big(\bQ\bu+\bd\big)
   =\boxed{\;\bK\bw+\bd_{\kappa}(t),\qquad
   \bd_{\kappa}=\bK\shift^{t}\bQ^{t}\bd\;}
\]
so the referential displacements differ by a rotation and a
spatially uniform translation. Taking $\Grad^{+}$ kills the translation
and returns \eqr{12.9}, as it must: a uniform displacement is a rigid
translation and strains nothing.
\end{solution}

\begin{problem}[Energy balance in the purely mechanical theory]
Show that in the purely mechanical theory
$\tfrac{\dl}{\dl t}E(S,t)=P(S,t)$, where $P$ is the power of the forces
acting on $S\subset B$ and
\[
   E=\int_{P}U\,\dl V+K(S,t),
\]
$K$ being the kinetic energy and $P=\kappa(S)$. The letter $P$ carries
two meanings in this statement, the power on the left and the reference
region under the integral sign; both are the book's, and they are told
apart by position.
\end{problem}

\begin{solution}
The mechanical energy balance \eqr{6.150} reads
\[
   \frac{\dl}{\dl t}K(S,t)+\mathcal{S}(S,t)=P(S,t),
\]
with the stress power $\mathcal{S}$ given in referential form by
\eqr{6.156}, $\mathcal{S}=\int_{P}\bP\cdot\dot\bF\,\dl V$. Everything
turns on identifying that integrand.

By \eqr{6.159}, $\bP\cdot\dot\bF=\bS\cdot\dot\bE$, the
Piola--Kirchhoff stress being power conjugate to the Lagrange strain.
By \eqr{12.32}, $\bS=U_{\bE}$, and the temperature being fixed, the
chain rule for $U(\bE;\bX)$ gives
\[
   \dot U=U_{\bE}\cdot\dot\bE=\bS\cdot\dot\bE=\bP\cdot\dot\bF .
\]
Therefore
\[
   \mathcal{S}(S,t)=\int_{P}\dot U\,\dl V
   =\frac{\dl}{\dl t}\int_{P}U\,\dl V ,
\]
the differentiation passing through the integral because $P$ is a fixed
region of $\kappa$: it is the reference placement of the fixed set $S$
of material points, and it does not move as the body deforms. This is
the one place where the referential form of the stress power is
essential; the spatial form would have a moving domain and the
interchange would need Reynolds' formula.

Substituting,
\[
   \frac{\dl}{\dl t}K+\frac{\dl}{\dl t}\int_{P}U\,\dl V=P(S,t),
   \qquad\text{that is}\qquad
   \boxed{\;\frac{\dl}{\dl t}E(S,t)=P(S,t)\;}
\]
with $E=\int_{P}U\,\dl V+K$.

\textbf{What is being said.} In a hyperelastic body at constant
temperature the working of the external forces goes entirely into
kinetic energy and stored strain energy, and none of it is lost. The
stress power, which Remark~\ref{rem:ch6stresspower} of Chapter 6
identified as the term with no counterpart in particle mechanics,
turns out here to be the rate of a potential. Chapter 9 had no such
statement, the viscous stress power being a dissipation; and
Chapter 11 needed the full Clausius--Duhem inequality to see that
$\bP=W_{\bF}$, which is what makes the potential exist.
\end{solution}

\begin{problem}[Residual stress in the linear theory]
Suppose the residually stressed undeformed configuration of an isotropic
material is in equilibrium without body force, and that the associated
traction vanishes on a portion of the boundary. Prove that the residual
stress vanishes everywhere in the body.
\end{problem}

\begin{solution}
This is the linearized twin of Problem 10 of Chapter 11, and it runs in
three steps of the same shape.

\textbf{Isotropy makes it hydrostatic.} By \eqr{12.50}$_{1}$ the
residual stress of an isotropic material is
\[
   \bS_{0}(\bX)=\alpha(\bX)\hI ,
\]
a scalar multiple of the referential identity at each point. Isotropy
permits no deviatoric part, because $U$ can depend on $\bE$ only through
its invariants, and at $\bE=\hOb$ the only tensor available from
$I_{1}=\hI\cdot\bE$ is $\hI$ itself.

\textbf{Equilibrium makes it uniform.} With no body force and no
acceleration, \eqr{12.71} reduces to $\Dvg\bS_{0}=\bze$. For a spherical
field this is a gradient,
\[
   \Dvg\big(\alpha\hI\big)=\Grad\alpha=\bze ,
\]
by the referential form of Problem 24 of Chapter 1, so $\alpha$ is
constant throughout the body, which is connected.

\textbf{The free patch makes it zero.} On the traction-free portion,
\eqr{12.64} gives $\bS_{0}\bN=\bze$ for the unit normal $\bN$ there,
that is $\alpha\bN=\bze$; and $|\bN|=1$, so $\alpha=0$. Being constant,
$\alpha$ vanishes identically and
\[
   \boxed{\;
\begin{aligned}
   &\bS_{0}=\hOb\ \text{ throughout the body},\\
   &\text{and by \eqr{12.54} and \eqr{12.55} so do }\bP
   \text{ and }\bT .
\end{aligned}\;}
\]

\textbf{Why the linear theory may assume \eqr{12.52}.} The result is the
justification for dropping $\bS_{0}$ from \eqr{12.51}. It is not a
theorem about all materials: an anisotropic body can hold residual
stress with a non-zero deviator, and the shrink-fitted ring and the
toughened glass plate are engineered to do exactly that. What the
problem shows is that isotropy plus equilibrium plus one free patch
leaves no room for it.
\end{solution}

\begin{problem}[The cofactor at leading order, and the two volume
formulas]
Show that
\[
   \bF^{*}-\shift\simeq\shift\big[(\tr\bG)\hI-\bG^{t}\big],
\]
and hence that $\bF^{*}\simeq\shift$ as claimed in the text. Use this to
confirm $\ip\bu\bn\,\dl a\simeq\ip\bw\bN\,\dl A$, and hence that the
right-hand sides of \eqr{12.26} and \eqr{12.27} agree at leading order.
\end{problem}

\begin{solution}
\textbf{The cofactor.} By \eqr{2.55}, $\bF^{*}=J\bF^{-t}$, so both
factors are needed to first order in $\bG$.

From \eqr{12.6}, $\bF=\shift(\hI+\bG)$, and inverting a product,
\[
   \bF^{-1}=\big(\hI+\bG\big)^{-1}\shift^{-1}
   \simeq\big(\hI-\bG\big)\shift^{t},
\]
the expansion $(\hI+\bG)^{-1}\simeq\hI-\bG$ being the geometric series
truncated after one term, legitimate because $|\bG|\ll1$, and
$\shift^{-1}=\shift^{t}$. Transposing,
\[
   \bF^{-t}\simeq\shift\big(\hI-\bG^{t}\big).
\]
From \eqr{12.25}, $J\simeq1+\tr\bG$. Multiplying and discarding the
product of two first-order terms,
\[
   \bF^{*}=J\bF^{-t}\simeq\big(1+\tr\bG\big)\shift
   \big(\hI-\bG^{t}\big)
   \simeq\shift\big[\hI+(\tr\bG)\hI-\bG^{t}\big],
\]
that is
\[
   \boxed{\;\bF^{*}-\shift\simeq\shift
   \big[(\tr\bG)\hI-\bG^{t}\big] .\;}
\]
The right-hand side has norm of order $|\bG|$, so $\bF^{*}\simeq\shift$
to leading order, as \eqr{12.55} and \eqr{12.62} assumed.

\textbf{The two volume formulas.} By Piola--Nanson \eqr{2.54} and the
estimate just proved,
\[
   \bn\,\dl a=\bF^{*}\bN\,\dl A\simeq\shift\bN\,\dl A .
\]
Taking the inner product with the spatial displacement $\bu$, and moving
the shifter across by the definition of the transpose,
\[
   \ip\bu\bn\,\dl a\simeq\ip\bu{\shift\bN}\,\dl A
   =\ip{\shift^{t}\bu}\bN\,\dl A
   =\boxed{\;\ip\bw\bN\,\dl A\;}
\]
by \eqr{12.4}. Integrating over the boundary,
\[
   \int_{\partial P_{t}}\ip\bu\bn\,\dl a
   \simeq\int_{\partial P}\ip\bw\bN\,\dl A ,
\]
which is exactly the equality of the right-hand sides of \eqr{12.27} and
\eqr{12.26}. Both compute the same $\Delta V$, one over the deformed
boundary and one over the reference boundary, and at leading order it
does not matter which.

\textbf{A check on the order of the errors.} Every discarded term above
was of order $|\bG|^{2}$, while $\Delta V$ itself is of order $|\bG|$
by \eqr{12.25}. So the relative error in $\Delta V$ is of order
$|\bG|$, which is the accuracy the whole chapter works to.
\end{solution}

\begin{problem}[The symmetries of the spatial stiffness]
Verify \eqr{12.59} and \eqr{12.60}.
\end{problem}

\begin{solution}
Everything follows from \eqr{12.58},
\[
   \mathbb{C}_{ijkl}=\dd_{iA}\dd_{jB}\dd_{kC}\dd_{lD}\,
   \mathbb{D}_{ABCD},
\]
in which the four shifter components are attached one to each index and
are the same for every term. Relabelling the referential summation
indices is therefore free, and the symmetries pass straight through.

\textbf{Major symmetry.} Interchange the index pairs $(ij)$ and $(kl)$
on the left, and correspondingly relabel $A\leftrightarrow C$,
$B\leftrightarrow D$ on the right:
\[
   \mathbb{C}_{klij}=\dd_{kA}\dd_{lB}\dd_{iC}\dd_{jD}\mathbb{D}_{ABCD}
   =\dd_{iC}\dd_{jD}\dd_{kA}\dd_{lB}\mathbb{D}_{CDAB}
   =\mathbb{C}_{ijkl},
\]
the middle step renaming the dummies $A\leftrightarrow C$,
$B\leftrightarrow D$ and the last using \eqr{12.38},
$\mathbb{D}_{CDAB}=\mathbb{D}_{ABCD}$. Hence
$\boxed{\mathbb{C}_{ijkl}=\mathbb{C}_{klij}}$.

\textbf{First minor symmetry.} Interchanging $i$ and $j$ interchanges
$A$ and $B$ after relabelling:
\[
   \mathbb{C}_{jikl}=\dd_{jA}\dd_{iB}\dd_{kC}\dd_{lD}\mathbb{D}_{ABCD}
   =\dd_{iA}\dd_{jB}\dd_{kC}\dd_{lD}\mathbb{D}_{BACD}
   =\mathbb{C}_{ijkl},
\]
by \eqr{12.40}. Hence $\boxed{\mathbb{C}_{ijkl}=\mathbb{C}_{jikl}}$.

\textbf{Second minor symmetry.} The same argument on the pair $(kl)$,
using \eqr{12.41}, gives
$\boxed{\mathbb{C}_{ijkl}=\mathbb{C}_{ijlk}}$.

\textbf{What each one is for.} The first minor symmetry says the stress
comes out symmetric, $T_{ij}=T_{ji}$, which is \eqr{6.117} and hence
Euler's second postulate. The second says only the symmetric part of
$\grad\bu$ enters, so $T_{ij}=\mathbb{C}_{ijkl}\varepsilon_{kl}$ and
rigid infinitesimal rotations produce no stress. The major symmetry is
the one that has no counterpart in the balance laws: it says the
stiffness derives from a potential, and it is what makes Problem 2
work. Counting independent constants, a general fourth-order tensor in
three dimensions has $81$ components; the two minor symmetries reduce
this to $36$, and the major symmetry to $21$.
\end{solution}

\begin{problem}[The strain energy as a quadratic form]
Establish that in the absence of residual stress the strain energy $U$
is approximated at leading order by
$\tfrac12\vek{\epsilon}\cdot\mathbb{D}[\vek{\epsilon}]$, or equivalently
by $\tfrac12\vek{\varepsilon}\cdot\mathbb{C}[\vek{\varepsilon}]$. Hence
show that $\bS$ and $\bT$ are approximated by the gradients
$U_{\vek{\epsilon}}$ and $U_{\vek{\varepsilon}}$.
\end{problem}

\begin{solution}
\textbf{The quadratic form.} Take the energy datum at the undeformed
state, $U(\hOb;\bX)=0$, which costs nothing since only differences of
energy have meaning. With $\bS_{0}=\hOb$, \eqr{12.33} reduces to
\[
   U(\bE;\bX)=\tfrac12\bE\cdot\mathbb{D}(\bX)[\bE]+o\big(|\bE|^{2}\big),
\]
and substituting $\bE\simeq\vek{\epsilon}$ from \eqr{12.8}, whose error
is of order $|\bG|^{2}$ and therefore contributes at order $|\bG|^{3}$
inside a quadratic form,
\[
   \boxed{\;U\simeq\tfrac12\,\vek{\epsilon}\cdot
   \mathbb{D}[\vek{\epsilon}]\;}
\]

\textbf{The spatial form.} Write both sides in components, using
\eqr{12.17}$_{2}$ in the form
$\epsilon_{AB}\simeq\dd_{iA}\dd_{jB}\varepsilon_{ij}$ and \eqr{12.58}:
\[
\begin{aligned}
   \vek{\epsilon}\cdot\mathbb{D}[\vek{\epsilon}]
   &=\epsilon_{AB}\epsilon_{CD}\mathbb{D}_{ABCD}\\
   &\simeq\dd_{iA}\dd_{jB}\varepsilon_{ij}\,
    \dd_{kC}\dd_{lD}\varepsilon_{kl}\,\mathbb{D}_{ABCD}
   =\varepsilon_{ij}\varepsilon_{kl}\mathbb{C}_{ijkl}
   =\vek{\varepsilon}\cdot\mathbb{C}[\vek{\varepsilon}],
\end{aligned}
\]
so
\[
   \boxed{\;U\simeq\tfrac12\,\vek{\varepsilon}\cdot
   \mathbb{C}[\vek{\varepsilon}]\;}
\]
The strain energy is the same number computed on either description, as
it must be, being a scalar.

\textbf{The stresses as gradients.} Differentiate the first boxed
expression with respect to $\vek{\epsilon}$. For a quadratic form with
major symmetry the gradient is twice the linear map, halved by the
factor $\tfrac12$: taking the rate,
\[
   \dot U\simeq\tfrac12\big\{\dot{\vek{\epsilon}}\cdot
   \mathbb{D}[\vek{\epsilon}]+\vek{\epsilon}\cdot
   \mathbb{D}[\dot{\vek{\epsilon}}]\big\}
   =\mathbb{D}[\vek{\epsilon}]\cdot\dot{\vek{\epsilon}},
\]
the two terms merging by \eqr{12.42}. Since $\dot{\vek{\epsilon}}$ is an
arbitrary symmetric tensor,
\[
   U_{\vek{\epsilon}}=\mathbb{D}[\vek{\epsilon}]\simeq\bS
\]
by \eqr{12.52} and \eqr{12.51}. The identical computation on the spatial
form gives $U_{\vek{\varepsilon}}=\mathbb{C}[\vek{\varepsilon}]\simeq\bT$
by \eqr{12.56}. So
\[
   \boxed{\;\bS\simeq U_{\vek{\epsilon}},\qquad
   \bT\simeq U_{\vek{\varepsilon}} .\;}
\]
Both stresses are gradients of the same potential, taken with respect to
the strain measure that lives in the same space. This is the linearized
form of $\bP=W_{\bF}$ from \eqr{11.16}$_{2}$, and it is why the major
symmetry \eqr{12.42} could not have been dispensed with.
\end{solution}

\begin{problem}[Isotropic moduli, and when the energy is positive
definite]
Find expressions for the components $\mathbb{C}_{ijkl}$ and
$\mathbb{D}_{ABCD}$ for isotropic solids. Linear elasticity always
assumes the strain-energy function positive definite, that is $U\ge0$
with equality only at zero strain. Show that for isotropic solids this
holds if and only if $\mu>0$ and $\kappa>0$, where
$\kappa=\lambda+\tfrac23\mu$.
\end{problem}

\begin{solution}
\textbf{The components.} By \eqr{12.50}$_{2}$,
$\mathbb{D}[\bE]=\lambda(\tr\bE)\hI+2\mu\bE$. In components this is
$\mathbb{D}_{ABCD}E_{CD}$, so we need a $\mathbb{D}_{ABCD}$ reproducing
$\lambda\dd_{AB}E_{CC}+2\mu E_{AB}$. Take
\[
   \boxed{\;\mathbb{D}_{ABCD}=\lambda\,\dd_{AB}\dd_{CD}
   +\mu\big(\dd_{AC}\dd_{BD}+\dd_{AD}\dd_{BC}\big)\;}
\]
and check by contracting with $E_{CD}$:
\[
   \mathbb{D}_{ABCD}E_{CD}
   =\lambda\dd_{AB}E_{CC}+\mu\big(E_{AB}+E_{BA}\big)
   =\lambda(\tr\bE)\dd_{AB}+2\mu E_{AB},
\]
using $E_{AB}=E_{BA}$. The expression has both minor symmetries by
inspection and the major symmetry because it is unchanged under
$(AB)\leftrightarrow(CD)$. By \eqr{12.58} the spatial components are the
same numbers,
\[
   \boxed{\;\mathbb{C}_{ijkl}=\lambda\,\dd_{ij}\dd_{kl}
   +\mu\big(\dd_{ik}\dd_{jl}+\dd_{il}\dd_{jk}\big)\;}
\]
the shifter components turning each referential Kronecker delta into a
spatial one.

\textbf{The quadratic form.} By Problem 6, and with $\bS_{0}=\hOb$,
\[
   2U=\vek{\epsilon}\cdot\mathbb{D}[\vek{\epsilon}]
   =\lambda\big(\tr\vek{\epsilon}\big)^{2}
   +2\mu\,\vek{\epsilon}\cdot\vek{\epsilon}.
\]
The two terms are not independent as they stand, because
$\vek{\epsilon}\cdot\vek{\epsilon}$ contains the trace. Split off the
spherical part: write
\[
   \vek{\epsilon}=\vek{\epsilon}'+\tfrac13\big(\tr\vek{\epsilon}\big)\hI,
   \qquad \tr\vek{\epsilon}'=0 ,
\]
the deviator of Chapter 1. The two pieces are orthogonal, since
$\vek{\epsilon}'\cdot\hI=\tr\vek{\epsilon}'=0$, so
\[
   \vek{\epsilon}\cdot\vek{\epsilon}
   =\vek{\epsilon}'\cdot\vek{\epsilon}'
   +\tfrac19\big(\tr\vek{\epsilon}\big)^{2}\,\hI\cdot\hI
   =|\vek{\epsilon}'|^{2}+\tfrac13\big(\tr\vek{\epsilon}\big)^{2},
\]
using $\hI\cdot\hI=\tr\hI=3$. Substituting,
\[
   2U=\Big(\lambda+\tfrac23\mu\Big)\big(\tr\vek{\epsilon}\big)^{2}
   +2\mu\,|\vek{\epsilon}'|^{2}
   =\boxed{\;\kappa\big(\tr\vek{\epsilon}\big)^{2}
   +2\mu\,|\vek{\epsilon}'|^{2}\;}
\]
with $\kappa=\lambda+\tfrac23\mu$.

\textbf{The criterion.} The two terms are now genuinely independent: the
scalar $\tr\vek{\epsilon}$ and the traceless tensor $\vek{\epsilon}'$
may be chosen separately, and together they determine
$\vek{\epsilon}$.

If $\mu>0$ and $\kappa>0$ then $2U$ is a sum of two non-negative terms,
so $U\ge0$; and $U=0$ forces both $\tr\vek{\epsilon}=0$ and
$\vek{\epsilon}'=\hOb$, hence $\vek{\epsilon}=\hOb$.

Conversely, suppose $U$ is positive definite. Taking
$\vek{\epsilon}=c\hI$ with $c\neq0$ gives $\vek{\epsilon}'=\hOb$ and
$\tr\vek{\epsilon}=3c$, so $2U=9\kappa c^{2}>0$ and hence $\kappa>0$.
Taking $\vek{\epsilon}$ any non-zero traceless tensor, for instance
$\hbE_{1}\otimes\hbE_{2}+\hbE_{2}\otimes\hbE_{1}$, gives
$\tr\vek{\epsilon}=0$ and $2U=2\mu|\vek{\epsilon}'|^{2}>0$, hence
$\mu>0$. So
\[
   U\ \text{positive definite}\iff \mu>0\ \text{ and }\ \kappa>0 .
\]

\textbf{Reading the two constants.} The decomposition separates the two
ways a body can be strained. A pure dilatation has
$\vek{\epsilon}'=\hOb$ and costs $\tfrac12\kappa(\tr\vek{\epsilon})^{2}$,
so $\kappa$ is the bulk modulus. A pure shear has $\tr\vek{\epsilon}=0$
and costs $\mu|\vek{\epsilon}'|^{2}$, so $\mu$ is the shear modulus,
the $G$ of \eqr{11.41}. Positive definiteness says both cost energy.
Note that $\lambda$ itself need not be positive: $\lambda=-\tfrac23\mu$
gives $\kappa=0$, which is excluded, but any $\lambda>-\tfrac23\mu$ is
admissible, and negative $\lambda$ corresponds to materials that
thicken laterally when compressed.
\end{solution}

\begin{problem}[The spatial equation of motion]
Show that Cauchy's equation of motion in an inertial frame is
approximated at leading order in the small displacement gradient by
\[
   \dvg\bT+\rho_{\kappa}\bb=\rho_{\kappa}\ddot\bu .
\]
\end{problem}

\begin{solution}
Start from the referential form \eqr{12.71} and push every term forward
with the shifter.

\textbf{The divergence.} By \eqr{12.55}, $T_{ij}\simeq\dd_{iA}\dd_{jB}
S_{AB}$. Differentiating with respect to $x_{j}$ and converting the
derivative by \eqr{12.14} in the form
$(\cdot)_{,j}\simeq(\cdot)_{,C}\dd_{jC}$,
\[
   T_{ij,j}\simeq\dd_{iA}\dd_{jB}S_{AB,C}\dd_{jC}
   =\dd_{iA}\dd_{BC}S_{AB,C}=\dd_{iA}S_{AB,B},
\]
using $\dd_{jB}\dd_{jC}=\dd_{BC}$, which is $\shift^{t}\shift=\hI$ in
components. Hence
\[
   \dvg\bT\simeq\shift\big(\Dvg\bS\big).
\]

\textbf{The other terms.} The body force satisfies
$\shift\bff=\shift\shift^{t}\bb=\bb$, and the acceleration
$\shift\ddot\bw=\shift\shift^{t}\ddot\bu=\ddot\bu$, both by
$\shift\shift^{t}=\I$.

\textbf{Assembling.} Apply $\shift$ to \eqr{12.71}:
\[
   \shift\big(\Dvg\bS\big)+\rho_{\kappa}\shift\bff
   =\rho_{\kappa}\shift\ddot\bw ,
\]
that is
\[
   \boxed{\;\dvg\bT+\rho_{\kappa}\bb=\rho_{\kappa}\ddot\bu\;}
\]
at leading order.

\textbf{Which density.} By \eqr{12.29}, $\rho\simeq\rho_{\kappa}$, so
the same equation may be written with $\rho$ in place of
$\rho_{\kappa}$: the difference is of first order in $|\bG|$ and
multiplies terms that are already of first order, so it is of second
order and lies below the accuracy of the theory. That is why the
literature is free to write either, and it is worth remembering that the
freedom is an approximation, not an identity. In the exact theory
\eqr{6.109} carries $\rho$ and \eqr{6.136} carries $\rho_{\kappa}$, and
the two equations are genuinely different.

For an isotropic material with constant moduli, substituting
\eqr{12.61} and using $\dvg[(\grad\bu)^{t}]=\grad(\dvg\bu)$ gives the
Navier equation
\[
   \mu\,\dvg\grad\bu+(\lambda+\mu)\grad(\dvg\bu)+\rho_{\kappa}\bb
   =\rho_{\kappa}\ddot\bu ,
\]
which is the equation almost every elasticity calculation begins from.
\end{solution}

\begin{problem}[The equation of motion for a second inertial observer]
Verify \eqr{12.73}. Show that $\ddot\bw\mapsto\ddot\bw^{+}=\bK\ddot\bw$
under a Galilean transformation, and hence that \eqr{12.76} holds if the
frame of $O^{+}$ is inertial.
\end{problem}

\begin{solution}
\textbf{\eqr{12.73}.} By Problem 1 and \eqr{12.72}, $\bS$ transforms
like $\bE$, so $\bS^{+}=\bK\bS\bK^{t}$, that is
$S^{+}_{AB}=K_{AC}K_{BD}S_{CD}$. The reference positions are related by
$\bX^{+}=\bK\bX+\bc^{+}$ with $\bc^{+}$ constant, so
$X^{+}_{B}=K_{BE}X_{E}+c^{+}_{B}$ and, $\bK$ being orthogonal,
$\partial X_{E}/\partial X^{+}_{B}=K_{BE}$. The chain rule then gives
\[
   \frac{\partial S^{+}_{AB}}{\partial X^{+}_{B}}
   =K_{AC}K_{BD}\,\frac{\partial S_{CD}}{\partial X_{E}}\,K_{BE}
   =K_{AC}\dd_{DE}\,S_{CD,E}
   =K_{AC}\,S_{CD,D},
\]
using $K_{BD}K_{BE}=\dd_{DE}$, which is $\bK^{t}\bK=\hI$. Hence
\[
   \boxed{\;\Dvg^{+}\bS^{+}=\bK\big(\Dvg\bS\big)\;}
\]

\textbf{The acceleration.} By Problem 1, $\bw^{+}=\bK\bw+\bd_{\kappa}$
with $\bd_{\kappa}=\bK\shift^{t}\bQ^{t}\bd$. Differentiating twice, and
noting that $\bK$ is constant in time by hypothesis,
\[
   \ddot\bw^{+}=\bK\ddot\bw+\ddot{\bd}_{\kappa} .
\]
Under a Galilean transformation \eqr{7.10} the rotation $\bQ$ is
constant and $\bc(t)=\bV t+\bc_{0}$ is linear in $t$. By \eqr{7.21} the
uniform part of the displacement is
$\bd(t)=\bc(t)-\bQ\shift\bK^{t}\bc^{+}$, which is therefore also linear
in $t$, so $\ddot{\bd}=\bze$; and $\bd_{\kappa}=\bK\shift^{t}\bQ^{t}\bd$
with all three tensors constant, so $\ddot{\bd}_{\kappa}=\bze$ too.
Hence
\[
   \boxed{\;\ddot\bw^{+}=\bK\ddot\bw\;}
\]

\textbf{Assembling.} Substituting into \eqr{12.74},
\[
   \Dvg^{+}\bS^{+}+\rho^{+}_{\kappa^{+}}\bff^{+}
   =\rho^{+}_{\kappa^{+}}\bK\ddot\bw
   =\rho^{+}_{\kappa^{+}}\ddot\bw^{+},
\]
which is \eqr{12.76}. So the second observer writes the same equation,
with the same moduli by \eqr{12.67}: linear elasticity is Galilean
invariant, and no experiment within it can distinguish one inertial
frame from another.
\end{solution}

\begin{problem}[A non-inertial second observer]
How is \eqr{12.76} modified if the frame of $O^{+}$ is not inertial?
\end{problem}

\begin{solution}
Only the acceleration changes; \eqr{12.73} and the constitutive
transformations \eqr{12.66} used nothing about the motion of the
frames.

\textbf{The extra terms.} For general $\bQ(t)$ and $\bc(t)$ the
accelerations are related by \eqr{7.9},
\[
   \ba^{+}-\bQ\ba=\ba_{\mathrm{f}},\qquad
   \ba_{\mathrm{f}}=\ddot\bc+2\bOm\big(\bv^{+}-\dot\bc\big)
   +\big(\dot\bOm-\bOm^{2}\big)\big(\bx^{+}-\bc\big),
\]
$\bOm=\dot\bQ\bQ^{t}$ being the spin of the relative motion. In the
linear theory $\ba=\ddot\bu$ and $\ba^{+}=\ddot\bu^{+}$, by
\eqr{12.68} applied in each frame, so
\[
   \ddot\bw^{+}=(\shift^{+})^{t}\ba^{+}
   =\bK\shift^{t}\bQ^{t}\big(\bQ\ba+\ba_{\mathrm{f}}\big)
   =\bK\shift^{t}\ba+\bK\shift^{t}\bQ^{t}\ba_{\mathrm{f}}
   =\bK\ddot\bw+(\shift^{+})^{t}\ba_{\mathrm{f}} ,
\]
using $(\shift^{+})^{t}=\bK\shift^{t}\bQ^{t}$ from \eqr{7.18} and
$\shift^{t}\ba=\ddot\bw$. Hence
$\bK\ddot\bw=\ddot\bw^{+}-(\shift^{+})^{t}\ba_{\mathrm{f}}$, and
substituting into \eqr{12.74},
\[
   \boxed{\;
\begin{aligned}
   \Dvg^{+}\bS^{+}+\rho^{+}_{\kappa^{+}}
   \big[\bff^{+}+\bff^{+}_{\mathrm{f}}\big]
   &=\rho^{+}_{\kappa^{+}}\ddot\bw^{+},\\
   \text{where}\quad
   \bff^{+}_{\mathrm{f}}&=(\shift^{+})^{t}\ba_{\mathrm{f}} .
\end{aligned}\;}
\]
The non-inertial observer writes the same equation with an additional
body force per unit mass, and Chapter 7 already named its three parts:
$\ddot\bc$ is the translational term, $2\bOm(\bv^{+}-\dot\bc)$ the
Coriolis term, and $(\dot\bOm-\bOm^{2})(\bx^{+}-\bc)$ the Euler and
centrifugal terms.

\begin{remark}[what the Coriolis term does to the
equation]\label{rem:ch12coriolis}
The apparent force is not merely an inconvenience of bookkeeping. Two
of its parts spoil the structure of \eqr{12.76}.

The Coriolis term contains $\bv^{+}$, hence $\dot\bu^{+}$, so the
equation acquires a first derivative in time. It ceases to be the wave
equation of \eqr{12.76} and becomes one with a gyroscopic term, which
is why the vibration of a rotating shaft cannot be analysed with the
non-rotating modes.

The centrifugal term contains $\bx^{+}$, hence $\bu^{+}$, so it adds a
term proportional to the displacement itself. It behaves like a
negative stiffness and can destabilize a structure that is stable at
rest, which is the mechanism of spin softening in turbine blades. Both
effects are linear in $\bu^{+}$ and belong properly to the linear
theory; nothing here required an additional approximation.
\end{remark}
\end{solution}
